\scanpage{23}
\mspage{14}
% archive sheet number: 23
% lines 1--3
Ces conditions At M1--At M3 \inserted{plus At M4 plus bas} ne concernent que \deleted{$\unsure{\Leftarrow|\leftarrow|\triangleleft}$} la relation $X\triangleleft F$. Allons un peu plus loin. Si $X,Y\in M$, on a tautologiquement
% lines 4--6
\[
\begin{aligned}
F_X\leq F_Y&\Longleftrightarrow\forall\deleted{\unsure{\underline{\Phi}\in\widetilde{\mathscr{F}}}}F\in\unsure{\mathscr{F}|\mathcal{F}},\ \deleted{\unsure{X}}Y\triangleleft F\Rightarrow X\triangleleft F\\
&\Longleftrightarrow\forall\unsure{\underline{\Phi}|\underline{\varphi}}\in\widetilde{\mathscr{F}},\ Y\in\underline{\Phi}\Rightarrow X\in\underline{\Phi}\\
&\Longleftrightarrow X\triangleleft F_Y\Longleftrightarrow X\in\widetilde{F}_Y.
\end{aligned}
\]
% lines 7--11
On écrit alors $X\leq Y$, c'est une relation de préordre, savoir la relation de préordre associée à la famille \unsure{préadmissible|quasi-admissible} $\widetilde{\mathscr{F}}$ dans $M$. Par At M1 b) c'est une relation d'\uline{ordre}. Ceci dit, on a une
% slanted
\margin{On écrit $\widetilde{X}=\widetilde{F}_X=\{Y\mid Y\triangleleft\deleted{\unsure{f}}F_X\}=\deleted{\unsure{\mathscr{F}}}M_{\leq X}$}

% lines 12--14, cancelled proposition
\deleted{\uline{Prop} Si $F\in\mathscr{F}$, $\inserted{\underline{\Phi}=}\widetilde{F}$ est une partie fermée de $M,\leq$.}
\deleted{Si $\unsure{C|E}\ \Psi$ est une partie fermée de $\underline{\Phi}$, alors $\Psi\in\widetilde{\mathscr{F}}$}

% lines 15--16
\uline{Prop} Soit $\underline{\Phi}\in\unsure{\mathscr{F}|\widetilde{\mathscr{F}}}$, alors toute partie fermée \inserted{$\Psi$} de $\underline{\Phi}$ est dans $\widetilde{\mathscr{F}}$.

% lines 17--23
Ceci \deleted{\unsure{renforce At}} est formellement plus fort que At M3 a), et lui est donc équivalent (compte tenu du reste). Si j'ai laissé At M3 a) sous forme apparemment plus faible, c'est pour lui donner un sens en termes de $(\mathscr{F},\leq)$, comme axiome d'existence de Sup majorés.

% lines 24--28
\deleted{\uline{At M4} Soit $F,G\in\mathscr{F}$. Pour que $\unsure{F}$}

Soit $F,G\in\mathscr{F}$, on pose
\[
\begin{aligned}
F\unsure{\gtrless|\bowtie}G&\overset{\text{déf}}{\Longleftrightarrow}\{F,G\}\text{ majoré dans }\mathscr{F}\Longleftrightarrow F\cup G\text{ existe}\\
&\Longleftrightarrow\widetilde{F}\cup\widetilde{G}\in\widetilde{\mathscr{F}}
\end{aligned}
\]
\[
X\gtrless Y\overset{\text{déf}}{\Longleftrightarrow}F_X\gtrless F_Y\quad\text{i.e.}\quad\widetilde{X}\cup\widetilde{Y}\in\widetilde{\mathscr{F}}\quad[\widetilde{X}=M_{\leq X}]
\]
\[
\text{\uline{At M4}}\quad F\gtrless G\Longleftrightarrow\forall X\in\widetilde{F},\ Y\in\widetilde{G},\text{ on a }X\gtrless Y
\]
% lines 29--30
\deleted{\unsure{At}} Pour les relations \unsure{données} $\ll$, $\unsure{\circ\mid|\circ|\mid\circ\mid}$ sur $M$, on pose
\[
\text{\uline{At M5}}\quad X\gtrless\deleted{\unsure{F|G}}Y\Longleftrightarrow\forall X'\in\widetilde{X}\setminus\widetilde{X}\cap\widetilde{Y},\ Y'\in\widetilde{X}\cap\widetilde{Y},\text{ on a }X'\mid\circ\mid Y'.
\]
